\chapter{Этаж~1. Чуть выше: облако и \texorpdfstring{$\varepsilon$}{epsilon}-оболочки}
\label{ch:floor1}

\section{Радиус этажа}
\label{sec:floor1-scale}

Этаж~1 --- не одна ячейка, а \emph{составной узел} на каузальной звезде \latGCKshort\ (\bcref{def:packing-types}).
Длина задаётся степенями числа соседей $\geomNeighborTwelve=12$:
\[
  \geomNeighborTwelve^{1}=12,\quad
  \geomNeighborTwelve^{2}=144,\quad
  \geomNeighborTwelve^{3}=1728,\quad
  \geomNeighborTwelve^{4}\approx 2.1\times 10^{4},\quad
  \geomNeighborTwelve^{5}\approx 2.5\times 10^{5}.
\]
Это окно $10^{3}\ldots 10^{5}$ планковских шагов~$\caStepSpace$ ---
\emph{до} комптоновской и боровской шкал ($10^{22}\ldots 10^{24}\cdot\caStepSpace$,
\bcref{prop:ladder-scale-steps}).

Здесь живут оболочки соседей вокруг ядра планкона и облако материи~$\rho_\Theta$,
аналог электронной плотности на атомном языке, но на~$\layerMicro$.

\section{Три слоя плотности}
\label{sec:floor1-density-layers}

На одном узле различают три объекта, которые на~$\layerMacro$ часто смешивают:

\begin{table}[htbp]
  \centering
  \caption{Слои плотности вокруг планкона.}
  \label{tab:floor1-density}
  \begin{tabular}{@{}lll@{}}
    \toprule
    поле & смысл на~$\layerMicro$ & аналог на~$\layerMacro$ \\
    \midrule
    $|z|^2$, $\rho_E$ & энергоплотность \emph{среды}; вакуум кипит везде & не $\rho_e$ \\
    $b\in\{0,1\}$ на ядре & занятость ячейки: дырка или планкон & жёсткий скелет узла \\
    $\rho_\Theta(x)$ & облако вокруг ядра: где $|\Delta\varphi|$, $|\zeta|\geq\Delta\varphi_{\min}$ & $\rho_e(r)$ \\
    \bottomrule
  \end{tabular}
\end{table}

Гейзенберг на голономии ($\Delta\varphi\geq\Delta\varphi_{\min}=\tfrac12$ рад,
\bcref{lem:51a}) не позволяет сжать фазовый дефект в бесконечно острый угол.
Поэтому планкон не может быть только битом «ячейка занята'':
вокруг ядра остаётся \emph{неустранимый ореол}.
Модуль упругости~$K_P$ (\bcref{def:kp}) удерживает ядро; Гейзенберг --- хвост.

На~$\Mlayer$ облако задаётся порогом, без подгоняемой гладкой функции:
\[
  \rho_\Theta(x)=\mathbf{1}\big[|\Delta\varphi_N(x)|\geq\Delta\varphi_{\min}\big].
\]
Гладкость появляется только после осреднения с биномиальным ядром (\bcref{def:soft-readout}) на~$\layerMacro$.

\section{Минимальный ореол: \texorpdfstring{$R_{\mathrm{dress}}=\caStepSpace$}{R=dl}}
\label{sec:floor1-dressing}

\begin{proposition}[Минимальный радиус ореола]
\label{prop:floor1-dress-radius}
У легчайшего заряженного состояния $Q=\pm1$ минимальный пространственный хвост
облака~$\rho_\Theta$ покрывает всю $\varepsilon$-звезду первых соседей:
\[
  R_{\mathrm{dress}}=R_{\min}=1\cdot\caStepSpace,\qquad N_{\mathrm{star}}=\geomNeighborTwelve.
\]
\end{proposition}

\begin{proof}
\pfrom{звезды из $\geomNeighborTwelve$ связей на узле \latGCKshort}
фазовый шаг за полный оборот по соседям равен
\[
  \Delta\varphi_{\mathrm{ring}}=\frac{2\pi}{\geomNeighborTwelve}.
\]
\pusing{числа секторов $N_\varphi=\lceil 2\pi/\Delta\varphi_{\min}\rceil=13$ и $\Delta\varphi_{\min}=\tfrac12$}
\[
  \frac{2\pi}{\geomNeighborTwelve}>\Delta\varphi_{\min}.
\]
\pfollows
вся $\varepsilon$-окрестность узла лежит выше гейзенберговского пола фазы.

\pfrom{локальности закона~$g$ (\Axref{ax:A1})}
за один такт~$\htick$ доступны только соседи $N(x)$;
вторая координационная сфера за один такт недостижима.
\phence минимальный устойчивый хвост облака $\rho_\Theta$ покрывает первую звезду, и
\[
  R_{\mathrm{dress}}=R_{\min}=1\cdot\caStepSpace.
\]
Возбуждение на нескольких оболочках с тем же $Q$ не является основным состоянием
и релаксирует к основному состоянию (\cref{sec:floor1-census,sec:floor1-excited-states}).
\end{proof}

Оболочки на этаже~1 --- это моды~$\rho_\Theta$ на координационных сферах,
не орбитали $r<\ell_P$ внутри пикселя (\cref{ch:floor0}).

\section{Координационные сферы и каталог мод}
\label{sec:floor1-shell-catalog}

На графе $\Lambda$ с каузальной звездой $N_{12}$ (FCC, \bcref{ch:carrier})
вокруг узла планкона различают \emph{координационные сферы} --- множества узлов
на фиксированном графовом расстоянии от ядра.

\begin{table}[htbp]
  \centering
  \caption{Сферы вокруг планкона на FCC ($N_{12}=12$).}
  \label{tab:floor1-shells}
  \begin{tabular}{@{}cll@{}}
    \toprule
    сфера $k$ & узлов & смысл для $\rho_\Theta$ \\
    \midrule
    $0$ & $1$ & ядро $b=1$ \\
    $1$ & $12$ & $\varepsilon$-звезда; $R_{\mathrm{dress}}=1\cdot\caStepSpace$ \\
    $2$ & $42$ & внешние моды возбуждения \\
    $3$ & $92$ & следующий ярус возбуждения \\
    \bottomrule
  \end{tabular}
\end{table}

Шар $k\le 3$ содержит $147$ узлов --- конечный носитель каталога на этом этаже,
без континуума внутри одной планковской ячейки.

\begin{proposition}[Основная мода и возбуждения на сферах $2$--$3$]
\label{prop:floor1-shell-modes}
У легчайшего $Q=\pm1$ устойчивое облако $\rho_\Theta$ занимает только сферу~$1$
(\cref{prop:floor1-dress-radius}).
Любая непустая активация сферы $k\ge 2$ при том же топозаряде --- \emph{возбуждённое состояние}
того же $Q=\pm1$, не новая устойчивая материя.
\end{proposition}

\begin{proof}
\pfrom{\cref{prop:floor1-dress-radius} и локальности A1}
основное состояние --- минимальный ореол на $\varepsilon$-звезде.
\pfrom{учёта классов B=0 на полосе этажа~1 (\cref{sec:floor1-census})}
единственная устойчивая материя с $B=0$ на этом радиусе --- легчайший $Q=\pm1$
с минимальным ореолом на $\varepsilon$-звезде;
возбуждения с тем же $Q$ релаксируют к этому основному состоянию с уходом энергии
в свободный фотонный сектор.
\phence активация сфер $2$ или $3$ не даёт нового устойчивого класса.
\end{proof}

\emph{Элементарные} моды на сферах $2$ и $3$: один узел с $\rho_\Theta=1$
при полной $\varepsilon$-звезде --- всего $42+92=134$ конфигурации.
Полный перечень всех составных паттернов на $42$ узлах сферы~$2$ не требуется:
они относятся к тем же возбуждённым состояниям, а не к отдельным частицам.

\subsection{Четыре канала кубооктаэдра и правила отбора между сферами}
\label{sec:floor1-four-channels}

Для любого узла $x\in\Lambda$ с упорядоченными модулями координат $a\ge b\ge c\ge 0$ ($a+b+c\equiv 0\pmod 2$)
графовое расстояние на \latGCKshort\ равно норме кубооктаэдра
\[
  k = \max\!\left(a,\;\frac{a+b+c}{2}\right).
\]
У каждого узла на сфере $k\ge 1$ его $\geomNeighborTwelve=12$ связей распадаются по сферам $k-1$, $k$ и $k+1$
($|\Delta k|\le 1$ за один такт~$\htick$ по A1):
\[
  d_{\mathrm{in}} + d_{\mathrm{same}} + d_{\mathrm{out}} = 12.
\]

\begin{theorem}[Четыре канала кубооктаэдра на координационных сферах ГЦК]
\label{thm:floor1-four-channels}
При любом $k\ge 1$ все $N_{\mathrm{shell}}(k)=10k^2+2$ узлов координационной сферы~$k$
исчерпываются ровно четырьмя классами $(d_{\mathrm{in}},d_{\mathrm{same}},d_{\mathrm{out}})$
по геометрическим элементам кубооктаэдра (пятого класса $d_{\mathrm{in}}\ge 5$ на \latGCKshort\ не существует):
\begin{enumerate}[noitemsep]
\item \textbf{Вершинный канал $V=(1,4,7)$} ($a=b=k,\ c=0$):
  $n_V(k)=12$ для всех $k\ge 1$ (продолжения $12$ лучей звезды; открыт с $k=1$);
\item \textbf{Рёберный канал $E=(2,5,5)$} ($a=k,\ b+c=k,\ c>0$):
  $n_E(k)=24(k-1)$ (внутренние точки $24$ рёбер; открыт с $k=2$);
\item \textbf{Квадратно-граневой канал $\square=(4,4,4)$} ($a=k,\ b+c<k$):
  $n_{\square}(k)=6(k-1)^2$ (внутренность $6$ квадратных граней; открыт с $k=2$);
\item \textbf{Треугольно-граневой канал $\triangle=(3,6,3)$} ($a<k,\ a+b+c=2k$):
  $n_{\triangle}(k)=4(k-1)(k-2)=8\binom{k-1}{2}$ (внутренность $8$ треугольных граней; открыт с $k=3$).
\end{enumerate}
При переходе $k-1\to k$ внешние связи $d_{\mathrm{out}}$ каждого узла образуют каскадную (верхнетреугольную) матрицу отбора:
\begin{align*}
  V\ (d_{\mathrm{out}}=7) &\longrightarrow 1\,V + 4\,E + 2\,\square + 0\,\triangle,\\
  E\ (d_{\mathrm{out}}=5) &\longrightarrow 0\,V + 2\,E + 2\,\square + 1\,\triangle,\\
  \square\ (d_{\mathrm{out}}=4) &\longrightarrow 0\,V + 0\,E + 4\,\square + 0\,\triangle,\\
  \triangle\ (d_{\mathrm{out}}=3) &\longrightarrow 0\,V + 0\,E + 0\,\square + 3\,\triangle.
\end{align*}
В частности, прямой однотактовый переход между вершинным лучом~$V$ и треугольной гранью~$\triangle$
запрещён ($d_{V\to\triangle}=0$): возбуждение $\triangle$-канала на сфере~$3$ идёт только двухтактовым мостом
через ребро сферы~$2$ ($V_{k=1}\xrightarrow{4} E_{k=2}\xrightarrow{1}\triangle_{k=3}$).
\end{theorem}

\section{Откуда берутся подоболочки \texorpdfstring{$s, p, d, f$}{s, p, d, f}: геометрия граней кубооктаэдра}
\label{sec:floor1-spdf-correspondence}

В континуальной атомной физике на~$\layerMacro$ четыре семейства подоболочек $s,p,d,f$
($\ell=0,1,2,3$ с размерностями $1,3,5,7$) получают решением дифференциального уравнения
для сферических гармоник $Y_{\ell m}(\theta,\varphi)$ на непрерывной сфере.
На решётке~\latGCKshort\ те же четыре семейства --- и причина того, почему их ровно четыре ---
вытекают непосредственно из геометрии четырёх элементов кубооктаэдра $(V,\square,\triangle,E)$
(\cref{thm:floor1-four-channels}) под действием кубической группы симметрии~$O_h$:

\begin{enumerate}[leftmargin=*]
\item \textbf{Сфера $k=1$ ($12$ вершин $V$, $(\pm 1,\pm 1,0)$) $\longrightarrow$ сферическая мода~$s$ ($\ell=0$, $1$ состояние).}
  На первой сфере открыт только вершинный канал~$V$.
  В основном состоянии (\cref{prop:floor1-dress-radius}) все $12$ вершин $\varepsilon$-звезды
  включены синфазно без узловых плоскостей (полносимметричное представление~$A_{1g}$),
  образуя решёточный прообраз $1s$-оболочки.

\item \textbf{Канал квадратных граней $\square$ (открывается на $k=2$: $6$ узлов $(\pm 2,0,0),(0,\pm 2,0),(0,0,\pm 2)$)
  $\longrightarrow$ три диполя $p_x,p_y,p_z$ ($\ell=1$, $3$ состояния) и два осевых квадруполя $d_{x^2-y^2},d_{z^2}$ ($e_g$, $2$ состояния).}
  На второй сфере $6$ центров квадратных граней лежат строго на трёх декартовых осях $\pm X,\pm Y,\pm Z$,
  образуя три пары противоположных квадратов (три антиподальные гантели вдоль осей).
  Пространство функций на этих $6$ осевых узлах раскладывается без остатка ($\mathrm{rank}=6$ из~$6$):
  \[
    6_{\square}
    = 1\,(s,\ A_{1g})
    \;\oplus\;
    3\,(p_x,p_y,p_z,\ T_{1u})
    \;\oplus\;
    2\,(d_{x^2-y^2},d_{z^2},\ E_g).
  \]
  Три $p$-орбитали существуют и ориентированы по осям $X,Y,Z$ потому, что у кубооктаэдра ровно $6$ квадратных граней
  ($3$ противоположные пары). Поскольку каждый такой узел опирается сразу на $d_{\mathrm{in}}=4$ активных узла
  $\varepsilon$-звезды планкона и имеет симметричный баланс связей $(4,4,4)$, квадратно-граневые моды
  удерживают возбуждение $\rho_\Theta=1$ в кипящем океане дольше всех остальных каналов сферы~$2$.

\item \textbf{Канал $12$ диагональных вершин $V$ ($(\pm k,\pm k,0)$ в плоскостях $XY,YZ,XZ$)
  $\longrightarrow$ три диагональных квадруполя $d_{xy},d_{yz},d_{xz}$ ($t_{2g}$, $3$ состояния) и полное замыкание $5$ мод $d$-оболочки ($\ell=2$).}
  На $6$ осевых узлах~$\square$ две координаты из трёх равны нулю, поэтому произведения $xy,yz,xz\equiv 0$
  и диагональные $d$-моды на~$\square$ тождественно обращаются в нуль ($\mathrm{rank}(d|_{\square})=2$).
  На $12$ вершинах кубооктаэдра~$V$ (по $4$ биссектрисных луча в каждой из трёх координатных плоскостей $XY,YZ,XZ$)
  открываются три крестообразные моды $d_{xy},d_{yz},d_{xz}$ ($T_{2g}$).
  Вместе с двумя осевыми модами $E_g$ на квадратах~$\square$ это даёт все $2+3=5$ состояний $d$-подоболочки
  ($\mathrm{rank}(d|_V)=5$ из~$5$), а пространство на $12$ вершинах~$V$ раскладывается без остатка ($\mathrm{rank}=12$ из~$12$):
  \[
    12_{V}
    = 1\,(s)
    \;\oplus\;
    3\,(p)
    \;\oplus\;
    5\,(d_{e_g\oplus t_{2g}})
    \;\oplus\;
    3\,\bigl(f_{x(y^2-z^2),\,y(z^2-x^2),\,z(x^2-y^2)}\bigr).
  \]
  Тем самым классическое расщепление пяти $d$-орбиталей на осевой дублет $e_g$ ($2$) и диагональный триплет $t_{2g}$ ($3$)
  заложено в саму геометрию вакуумного кубооктаэдра \latGCKshort: $e_g$ сидит на $6$ квадратных гранях~$\square$,
  а $t_{2g}$ --- на $12$ диагональных вершинах~$V$.

\item \textbf{Канал треугольных граней $\triangle$ (открывается на $k=3$: $8$ узлов $(\pm 2,\pm 2,\pm 2)$) и $24$ ребра $E$
  $\longrightarrow$ октуполь $f_{xyz}$ и полное замыкание $7$ мод $f$-оболочки ($\ell=3$).}
  Ни на $6$ центрах квадратных граней~$\square$, ни на $12$ вершинах~$V$ нет точек, где все три координаты
  одновременно отличны от нуля ($xyz\equiv 0$), поэтому октупольная мода $f_{xyz}$ на~$V$ и~$\square$
  невозможна ($\mathrm{rank}(f|_V)=6$ из~$7$).
  Только на третьей сфере ($k=3$) открывается четвёртый канал кубооктаэдра --- $8$ центров треугольных граней~$\triangle$
  с координатами $(\pm 2,\pm 2,\pm 2)$ в восьми октантах куба.
  Пространство на этих $8$ узлах раскладывается без остатка ($\mathrm{rank}=8$ из~$8$, бином $1+3+3+1=2^3$):
  \[
    8_{\triangle}
    = 1\,(s)
    \;\oplus\;
    3\,(p_x,p_y,p_z)
    \;\oplus\;
    3\,(d_{xy},d_{yz},d_{xz})
    \;\oplus\;
    1\,(f_{xyz},\ A_{2u}),
  \]
  а на $24$ внутренних точках рёбер~$E$ ($(\pm 2,\pm 1,\pm 1)$, где $x,y,z\neq 0$) одновременно линейно независимы
  все $7$ из $7$ мод $f$-подоболочки ($\mathrm{rank}(f|_E)=7$).
\end{enumerate}

\begin{table}[htbp]
  \centering
  \caption{Соответствие четырёх каналов кубооктаэдра ГЦК и подоболочек $s,p,d,f$.}
  \label{tab:floor1-spdf-channels}
  \footnotesize
  \setlength{\tabcolsep}{3.5pt}
  \begin{tabular}{@{}llllp{0.34\linewidth}@{}}
    \toprule
    элемент & открытие & носитель & разложение (ранг) & какие моды $(\ell=0,1,2,3)$ несёт \\
    \midrule
    $V$ ($12$ вершин) & $k=1$ & $12$: $(\pm k,\pm k,0)$ & $12 = 1_s+3_p+5_d+3_f$ &
      сфера $s$ ($1/1$); диагонали $d_{t_{2g}}$ ($3$) и вся $d$ ($5/5$); без $f_{xyz}$ ($6/7$) \\
    $\square$ ($6$ квадратов) & $k=2$ & $6$: $(\pm 2,0,0)$ & $6 = 1_s+3_p+2_{d,e_g}$ &
      гантели $p_x,p_y,p_z$ ($3/3$) по осям $\pm X,\pm Y,\pm Z$ + дублет $d_{e_g}$ ($2/5$) \\
    $\triangle$ ($8$ треуг.) & $k=3$ & $8$: $(\pm 2,\pm 2,\pm 2)$ & $8 = 1_s+3_p+3_{d,t_{2g}}+1_{f_{xyz}}$ &
      октуполь $f_{xyz}$ ($A_{2u}$) в $8$ октантах куба ($xyz\neq 0$) \\
    $E$ ($24$ ребра) & $k=2$ & $24$: $(\pm 2,\pm 1,\pm 1)$ & ранг $16$ на $s\oplus p\oplus d\oplus f$ &
      полное замыкание всех $7$ из $7$ мод $f$-подоболочки \\
    \bottomrule
  \end{tabular}
\end{table}

\section{Кто устойчив на этаже~1}
\label{sec:floor1-census}

На радиусах $\geomNeighborTwelve^{3}\ldots \geomNeighborTwelve^{4}$ рассматривают классы объектов с $B=0$
(нейтральный сектор по магнитному и топологическому бюджету узла).
Критерий устойчивости: нет пути вниз по~$g$ с теми же инвариантами и меньшей энергией.

\begin{table}[htbp]
  \centering
  \caption{Классы на лептонном окне этажа~1.}
  \label{tab:floor1-census}
  \footnotesize
  \setlength{\tabcolsep}{4pt}
  \begin{tabular}{@{}llp{0.42\linewidth}@{}}
    \toprule
    класс & устойчив? & комментарий \\
    \midrule
    фон A5 (кипящий вакуум) & нет & не объект, среда \\
    свободный фотон & да & уже вакуум; не сгусток радиуса $\geomNeighborTwelve^{k}$ \\
    основное $Q=\pm1$ + минимальный ореол & \textbf{да} & ядро этажа~0 + $\rho_\Theta$; \cref{prop:floor1-dress-radius} \\
    возбуждённое состояние (тот же $Q$) & нет & релаксация к основному + фотон \\
    нейтральный многоячейочный сгусток ($Q=0$) & нет & без топозаряда; распадается в фон, не материя \\
    пара $\pm$ на этом $R$ & нет & аннигиляция \\
    нейтральный $L$-сектор & возможно & вне полосы $\geomNeighborTwelve^{3\ldots 4}$ \\
    \bottomrule
  \end{tabular}
\end{table}

В окне $\geomNeighborTwelve^{3}\ldots \geomNeighborTwelve^{5}$ встречаются также
\emph{нейтральные многоячейочные сгустки} без топозаряда: они не защищены топологией
и не остаются как отдельные частицы --- только кратковременные узоры на фоне A5.
Единственная устойчивая \emph{материя} $B=0$ на этом радиусе --- легчайший $Q=\pm1$
с минимальным ореолом на $\varepsilon$-звезде.

\section{Возбуждённые состояния и релаксация}
\label{sec:floor1-excited-states}

\emph{Возбуждённое состояние} здесь --- тот же топологический заряд $Q=\pm1$, что и в основном,
но облако $\rho_\Theta$ занято ещё на координационных сферах $k\ge 2$, а не только на $\varepsilon$-звезде
(\cref{prop:floor1-dress-radius}).
По закону~$g$ существует нисходящий путь к основному состоянию с испусканием
электромагнитного кванта (свободный фотонный сектор, \cref{sec:floor1-census})
при сохранении аддитивных инвариантов узла;
поэтому такая конфигурация не является устойчивой материей на этаже~1.
Это та же логика, что \emph{релаксация} возбуждённого уровня в квантовой механике:
спуск в основное состояние и уход лишней энергии в излучение, а не распад «другой частицы''.

На макроязыке~$\layerMacro$ для многих таких событий ту же картину описывают
экспоненциальным распадом и шириной линии~$\Gamma$;
на~$\layerMicro$ длительность релаксации задаётся числом тактов
$n_{\mathrm{ticks}}\in\mathbb{N}\cdot\htick$, а не непрерывным временем жизни из PDG.
Каузальная нижняя граница для релаксации со сферы $k\ge 2$ равна $n_{\mathrm{ticks}}^{\min}(k)=k-1$ (из A1).
На заполненной решётке под A5 (кипящий океан на всех~$\hv$) одноузловое возбуждение с полуоборотом Паули
($N_{\mathrm{ring}}/2=256$ тиков при шаге Гейзенберга $\Delta\varphi_{\mathrm{disc}}=41$ тик)
отдаёт избыток во внешнюю ванну через долю связей $d_{\mathrm{out}}/\geomNeighborTwelve$,
что задаёт канальный потолок удержания (\cref{thm:floor1-four-channels}):
\begin{equation}
\label{eq:floor1-channel-ceiling}
  n_{\mathrm{ticks}}^{\max}(d_{\mathrm{out}})
  = \left\lceil\frac{\geomNeighborTwelve}{d_{\mathrm{out}}}\cdot\frac{N_{\mathrm{ring}}/2}{\Delta\varphi_{\mathrm{disc}}}\right\rceil
  = \left\lceil\frac{12}{d_{\mathrm{out}}}\cdot\frac{256}{41}\right\rceil
  =
  \begin{cases}
    11, & V\ (d_{\mathrm{out}}=7),\\
    15, & E\ (d_{\mathrm{out}}=5),\\
    19, & \square\ (d_{\mathrm{out}}=4),\\
    25, & \triangle\ (d_{\mathrm{out}}=3).
  \end{cases}
\end{equation}
На сфере~$2$ квадратно-граневые моды~$\square$ ($d_{\mathrm{in}}=4$, $d_{\mathrm{out}}/d_{\mathrm{in}}=1$)
опираются сразу на четыре активные вершины $\varepsilon$-звезды планкона и потому удерживают $\rho_\Theta=1$
до своего потолка ($\sim 19$ тактов), значительно дольше рёберных ($E$) и вершинных ($V$) мод.

\section{Связь с электроном}
\label{sec:floor1-electron-bridge}

Лабораторный электрон --- не объект этажа~0 и не весь этаж~1 целиком.
Схема:
\begin{enumerate}[noitemsep]
\item этаж~0: планкон ($b=1$, внутренний спектр, \cref{ch:floor0});
\item этаж~1: неустранимый ореол~$\rho_\Theta$ на $\varepsilon$-звезде;
\item дальше: составной узел (этаж~2), IR-лестница, атом на~$\layerMacro$
  (\cref{ch:floor2}, \bcref{ch:si-sm}).
\end{enumerate}

Масса $\massElectron$, заряд $\chargeE$ и комптоновская длина --- на оси $k\approx 8$,
не на $\geomNeighborTwelve^{m}$.
Путать $R_{\mathrm{dress}}=\caStepSpace$ с $\bar\lambda_C$ или $a_0$ --- частая ошибка;
\cref{ch:floors-preface}.

\section{Статус этажа~1}
\label{sec:floor1-open}

\paragraph{Закрыто.}
Окно $N_{12}^{3\ldots5}\cdot\caStepSpace$, учёт классов полосы (\cref{sec:floor1-census}),
$R_{\mathrm{dress}}=1\cdot\caStepSpace$,
форма $\rho_\Theta=\mathbf{1}[|\Delta\varphi|\ge\Delta\varphi_{\min}]$,
геометрия сфер $12/42/92$ и классификация мод (\cref{sec:floor1-shell-catalog,prop:floor1-shell-modes}).

\begin{table}[htbp]
  \centering
  \caption{Закрытые вопросы этажа~1 (оболочки и учёт).}
  \label{tab:floor1-closed}
  \begin{tabular}{@{}p{0.38\linewidth}p{0.52\linewidth}@{}}
    \toprule
    вопрос & итог \\
    \midrule
    радиус ореола легчайшего $Q=\pm1$ &
    $R_{\mathrm{dress}}=1\cdot\caStepSpace$, сфера $1$ (\cref{prop:floor1-dress-radius}) \\
    форма облака $\rho_\Theta$ &
    индикатор Гейзенберг-пола (\cref{sec:floor1-density-layers}) \\
    каталог на сферах $2$--$3$ &
    $42+92$ узлов; $134$ элементарных мод; возбуждённые состояния при активации (\cref{prop:floor1-shell-modes}) \\
    четыре канала отбора на сферах $k\ge 1$ &
    $V=(1,4,7)$, $E=(2,5,5)$, $\square=(4,4,4)$, $\triangle=(3,6,3)$; сумма $10k^2+2$; каскад и запрет $V\leftrightarrow\triangle$ за $1$ такт (\cref{thm:floor1-four-channels}) \\
    устойчивая материя $B=0$ на полосе &
    только основное $Q=\pm1$ (\cref{sec:floor1-census}) \\
    \bottomrule
  \end{tabular}
\end{table}

\paragraph{Смежные вопросы и связь с этажом~2.}

Для любого составного паттерна $S\subset S_2$ из $2^{42}$ конфигураций на второй координационной сфере скорость ухода фазового возбуждения в тепловую ванну определяется средним внешним ветвлением активных каналов $\bar d_{\mathrm{out}}(S) = \frac{1}{|S|}\sum_{i\in S} d_{\mathrm{out}}(i) \in [3, 7]$, поэтому время удержания каждого паттерна детерминированно заключено в интервале четырех одноузловых потолков~\eqref{eq:floor1-channel-ceiling}:
\[
  11 = n_{\mathrm{ticks}}^{\max}(V) \;\le\; n_{\mathrm{ticks}}(S) \;\le\; n_{\mathrm{ticks}}^{\max}(\triangle) = 25\ \mathrm{тактов}.
\]
Что касается отношения масс поколений $m_\mu/\massElectron \approx 206{,}7703$ и $m_\tau/\massElectron \approx 3477{,}47$, оно относится не к геометрическому окну одной оболочки этажа~1, а к трёхпоколенческому стеку на инфракрасной лестнице и полностью выведено в следующей главе (\cref{ch:floor2,sec:floor2-quark-masses-natural}) через фазовый сдвиг двухкомпонентного спинора $\delta_0 = 2/\lfloor\caBitBudget\rfloor = 2/9$.

\begin{table}[htbp]
  \centering
  \caption{Замыкание смежных вопросов этажа~1 и моста к этажу~2.}
  \label{tab:floor1-open}
  \begin{tabular}{@{}p{0.38\linewidth}p{0.52\linewidth}@{}}
    \toprule
    вопрос & итог \\
    \midrule
    распределение $n_{\mathrm{ticks}}$ для составных паттернов на сфере~$2$ ($2^{42}$) &
    \textbf{закрыто}: одноузловые каналы $(V,E,\square,\triangle)$ и каскад выведены (\cref{thm:floor1-four-channels}); любой составной паттерн ограничен вилкой $11\le n_{\mathrm{ticks}}(S)\le 25$ по среднему ветвлению $\bar d_{\mathrm{out}}(S)\in[3,7]$ (\cref{sec:floor1-excited-states}) \\
    отношение масс $m_\mu/\massElectron$ &
    \textbf{закрыто на этаже~2}: выведено без свободных параметров из $d=3$ стеков поколений с фазой $\delta_0=2/\lfloor\caBitBudget\rfloor=2/9$, $m_\mu/\massElectron\approx 206{,}7703$ ($+9{,}83\ \mathrm{ppm}$, \cref{sec:floor2-quark-masses-natural}) \\
    \bottomrule
  \end{tabular}
\end{table}

Следующий шаг вверх --- \cref{ch:floor2}: конфайнмент и кварковые фокусы
на радиусе порядка $10^{15}\cdot\caStepSpace$.
